% Generated draft: requires mathematical and source review.
\documentclass[11pt]{article}
\usepackage[margin=25mm]{geometry}
\usepackage{fontspec}
\setmainfont{Latin Modern Roman}
\usepackage{amsmath,amssymb,mathtools,unicode-math}
\setmathfont{Latin Modern Math}
\usepackage{enumitem,needspace,xurl,hyperref}
\setlength{\parindent}{0pt}
\setlength{\parskip}{5pt}
\setlength{\emergencystretch}{3em}
\newcommand{\N}{\mathbb{N}}
\newcommand{\Z}{\mathbb{Z}}
\newcommand{\Q}{\mathbb{Q}}
\newcommand{\R}{\mathbb{R}}
\newcommand{\C}{\mathbb{C}}
\newcommand{\F}{\mathbb{F}}
\newcommand{\A}{\mathbb{A}}
\newcommand{\PP}{\mathbb{P}}
\newcommand{\E}{\mathbb{E}}
\newcommand{\eps}{\varepsilon}
\newcommand{\dd}{\,\mathrm{d}}
\begin{document}
\section*{UnsolvedMath ID: 20000633}
\textbf{Source problem code:} catalogue entry
\subsection{Definitions and mathematical statement}
Maurice Rojas's source question is: what is the probability that a random Viro diagram contains no sphere? The record specifies no probability space, ambient dimension, support, subdivision, sign law, gluing convention, or meaning of ``sphere''. Consequently no single numerical probability is determined by that question as supplied. The following are separate models and reported results from the accompanying research aid.

In the aid's globally orthant-glued planar toric model, let $A\subset\Z^2$ be the nonempty set of used triangulation vertices, with independent signs satisfying $\Pr(\sigma_a=+1)=p$, $0\leq p\leq1$. The planar event means absence of a component homeomorphic to $S^1$ in the compact glued diagram. Set $A\bmod2=\{a+2\Z^2:a\in A\}$. The reported conditional probability is
\[
 \Pr(\hbox{no }S^1\hbox{ component}\mid A,\hbox{subdivision})
 =\mathbf1_{\{|A\bmod2|=1\}}\bigl[p^{|A|}+(1-p)^{|A|}\bigr].
\]
If the support and subdivision are random, with these conditional independent sign laws, the unconditional probability is the expectation of this same parity-and-cardinality expression. The aid reports probability zero for every primitive planar triangulation, where primitive means its lattice triangles are unimodular. The precise toric gluing and admissibility hypotheses are not supplied, so this formula remains qualified by that particular advertised model.

A different, unglued model is a positive-orthant polygonal fan with one interior vertex and $m$ boundary vertices, again with independent Bernoulli signs. Its reported probability of no local closed loop is
\[
 1-\bigl[p(1-p)^m+(1-p)p^m\bigr].
\]
This concerns a local loop inside the fan, and is not the preceding globally glued event.

For the abstract Bernoulli Viro-subcomplex model, a precise usual interpretation assigns values $\pm1$ to vertices of a triangulation, extends linearly on each simplex, and takes the zero level set. A vertex with spherical link whose sign differs from every adjacent sign yields an isolated component homeomorphic to $S^{n-1}$ in dimension $n$. Under independent fair signs, the aid reports the bound
\[
 \Pr(\hbox{no }S^{n-1}\hbox{ component})
 \leq (1-2^{-\Delta})^{N/(\Delta^2+1)},
\]
where $N$ counts eligible vertices with spherical links and $\Delta\geq1$ bounds the degree of the triangulation's adjacency graph. Application requires the local link and degree hypotheses; arbitrary boundary vertices need not qualify. Finally, for a connected closed triangulated surface with $N$ vertices in this abstract model, it reports the exact probability
\[
 \Pr(\hbox{no circle component})=p^N+(1-p)^N.
\]
These model-specific statements do not supply the missing definition of the original random diagram or identify which sphere dimension the original question intended.
\subsection{Short English statement}
The original question has no specified random model. The supplied aid gives separate formulas for a glued planar toric patchwork, a local planar fan, and an abstract signed triangulation. Their events and probability spaces differ, so none is a universal answer to the source question.
\subsection{Sources}
\begin{itemize}
\item Local UnsolvedMath registry, permanent ID $20000633$, source code \texttt{catalogue entry}, supplied complete source snapshot and research summary; canonical local record \texttt{aim-analysis-notes.json}, \texttt{notes[164]}.
\item Maurice Rojas, Question 8, AIM workshop \emph{Random analytic functions}, compiled by Swaminathan Sethuraman, April 17, 2006, as identified in the supplied verification note: \url{https://aimath.org/WWN/randomzeros/arcc1.pdf}. This is the supplied citation; no PDF was retrieved.
\item The original literature field is empty. No authoritative bibliographic reference or complete model definition for the displayed probability assertions is supplied.
\item Research attempt 1: partial result; problem status unknown at the run, registry status partially solved; supplied rigor label ``proved'', novelty confidence ``low''. The machine-generated aid is not peer reviewed and requires independent verification before citation.
\end{itemize}
\end{document}
